\section{The finite classifications}\label{sec:computation}

The path family explains why the inequality $\rho\le n+2$ eventually
fails for stable tangent bundles.  We now determine where the first
failure occurs in the classifications of smooth toric Fano varieties.
Throughout this section stability refers to the polarization $-K_X$.

\begin{theorem}\label{census:lowdim}
Let $X$ be a smooth toric Fano variety of dimension $n\le7$.
If $T_X$ is stable, then $\rho(X)\le n+2$.
For each $n=4,5,6,7$ this bound is attained, by respectively
$3,1,10,15$ isomorphism classes.
\end{theorem}

\begin{theorem}\label{census:eight}
The largest Picard number of a smooth toric Fano eightfold with stable
tangent bundle is $11$.  Exactly one isomorphism class attains it:
the path variety $X_2$ of Example~\ref{graph:path}.
\end{theorem}

\subsection{An independent test using bases}

For $S\subseteq\mathcal R$, put $\delta(S)=\sum_{\rho\in S}\delta_\rho$.
Lemma~\ref{lem:klyachko} is equivalent to
\begin{equation}\label{eq:subsetstability}
 \delta(S)<r(S)\quad
 \text{for every nonempty proper }S\subset\mathcal R,
\end{equation}
with weak inequalities for semistability. Indeed, if $r(S)<n$, apply
the subspace criterion to the span of $S$. If $r(S)=n$, positivity of
the weights gives $\delta(S)<\delta(\mathcal R)=n$. Conversely, for a
proper subspace $V$ containing rays, use the set of all rays in $V$;
its rank is at most $\dim V$.

A basis of the ray matroid is a subset $B\subset\mathcal R$ for which
$\{u_\rho:\rho\in B\}$ is a basis of $N_\CC$. Write $\mathbf1_B$ for
its incidence vector. The convex hull of these vectors is the
\emph{base polytope} of the ray matroid.

\begin{lemma}\label{lem:basecertificate}
Suppose the normalized degree vector admits a decomposition
\[
 \delta=\sum_{k=1}^{\ell}\lambda_k\mathbf1_{B_k},
 \qquad \lambda_k>0,\qquad\sum_{k=1}^{\ell}\lambda_k=1,
\]
where each $B_k$ is a basis of the ray matroid. Then $T_X$ is
semistable. Form a directed graph on $\mathcal R$ with an edge
$e\longrightarrow e'$ whenever, for some $k$, one has
$e\notin B_k$, $e'\in B_k$, and $(B_k\setminus\{e'\})\cup\{e\}$ is
also a basis. If this graph is strongly connected, then $T_X$ is stable.
\end{lemma}

\begin{proof}
For every $S\subseteq\mathcal R$,
\[
 \delta(S)=\sum_k\lambda_k|B_k\cap S|\le r(S),
\]
because $B_k\cap S$ is independent. Equality holds precisely when
$|B_k\cap S|=r(S)$ for every $k$.
If a nonempty proper $S$ gives equality and the directed graph is
strongly connected, an edge $e\longrightarrow e'$ leaves $S$.
For its corresponding basis $B_k$, replacing $e'\notin S$ by $e\in S$
produces a basis containing $r(S)+1$ members of $S$, a contradiction.
Thus \eqref{eq:subsetstability} holds. This argument does not use
Proposition~\ref{prop:symmetry}.
\end{proof}

\subsection{Coverage and exact slope comparisons}

We use the complete classification of smooth Fano polytopes obtained
by \O bro's algorithm \cite{Obro2007}, with the low-dimensional data
from polyDB and the seven- and eight-dimensional lists distributed
by Paffenholz \cite{polyDB}.  A smooth Fano polytope is the convex hull
of the primitive ray generators, and its facets determine the maximal
cones.  The facet normals of the dual moment polytope thus give the
same ray configuration.  This distinction matters because the two
distributions use different polytope conventions.

The computations through dimension six are those of
\cite{Wuebben2026}.  In dimension four, the classification by tangent
bundle stability was already obtained by Reynolds
\cite[Appendix~A]{Reynolds2023}; our four-dimensional calculation is an
independent verification, rather than a new classification.
The additional coverage used here is precisely
\[
 \begin{array}{c|c|r}
 n&\text{Picard numbers examined}&\text{number of varieties}\\ \hline
 7&\rho\ge9&1277\\
 8&\rho\ge11&1582.
 \end{array}
\]
Thus all sevenfolds at or above $\rho=n+2$ are covered, whereas in
dimension eight the complete calculation starts at $\rho=n+3$.
The eight-dimensional stratum $\rho=10$ has not been classified here
by stability.  No stability classification in dimension nine is used.

For each ray configuration, the intersection numbers $d_v$ are
computed as normalized facet volumes of the moment polytope.
Candidate facets and vertices obtained numerically are checked by
rational linear algebra before evaluating the volumes.  The identities
\[
 \sum_vd_v=D,\qquad \sum_vd_vv=0
\]
provide consistency checks, while each final slope comparison is a
comparison of rational numbers.  To disprove stability it suffices to
exhibit a proper ray-spanned subspace $W$ for which
\begin{equation}\label{census:witness}
 \frac{\sum_{v\in W}d_v}{\dim W}\ge\frac Dn.
\end{equation}
The search first examines subspaces spanned by at most four rays.
If it finds no such subspace, all ray-spanned subspaces are examined.
Consequently failure to find a small-dimensional subspace is never
used as evidence of stability.

There are two distinct kinds of non-stability record.  A strict
inequality in \eqref{census:witness} proves instability.  An equality
excludes stability but does not exclude the existence of a different
subspace with strictly larger slope.  We therefore retain the term
\emph{equality witness} for the latter records, without classifying
all of them as semistable.

\begin{proof}[Proof of Theorem~\ref{census:lowdim}]
The bound in dimensions at most six follows from the complete
low-dimensional calculation cited above and, in dimensions one and
two, from the classifications of smooth toric Fano curves and surfaces.
The numbers of varieties with $\rho\ge n+3$ in dimensions $4,5,6$
are $2,7,52$; none has stable tangent bundle.
Of the $1277$ sevenfolds examined, $1013$ have $\rho=9$ and $264$
have $\rho\ge10$.  Among the latter, $100$ admit a strict witness
and $164$ have a recorded equality witness.  Hence none is stable.
Among the former, exactly $15$ are stable.  Together with the stable
counts $3,1,10$ at $\rho=n+2$ in dimensions $4,5,6$, this proves both
assertions.
\end{proof}

\begin{table}[ht]
\centering
\begin{adjustbox}{max width=\textwidth}
\begin{tabular}{rrrrr}
\toprule
$\rho$ & total & stable & unstable & equality witness\\
\midrule
\csname @@input\endcsname tables/screen8_counts.tex
\bottomrule
\end{tabular}
\end{adjustbox}
\caption{All smooth toric Fano eightfolds with $\rho\ge11$.
The final column excludes stability but does not assert semistability.}
\label{census:screen8}
\end{table}

\begin{proof}[Proof of Theorem~\ref{census:eight}]
Table~\ref{census:screen8} accounts for all $1582$ varieties with
$\rho\ge11$: one is stable, $762$ have a strict destabilizing witness,
and $819$ have an equality witness.  The stable variety has $\rho=11$
and identifier \texttt{P8D.smooth\_fano\_8d.gz.389596} in Paffenholz's
list.  Its anticanonical degree is $D=5753440$.

Here is an explicit identification with $X_2$.  Order the coordinates
of its lattice as
\[
 (r_0,q_0,r_1,q_1,f^0_1,f^1_1,f^1_2,f^2_1),
\]
with the orientations in Example~\ref{graph:path}.  In the database
coordinates the lattice map is
\begin{equation}\label{census:identification}
 A=\begin{pmatrix}
 0&0&0&0&0&-1&0&0\\
 -1&0&0&0&0&0&0&0\\
 0&0&1&-1&0&0&0&0\\
 0&0&0&0&0&0&0&-1\\
 0&0&0&0&-1&0&0&0\\
 0&0&0&0&0&0&-1&0\\
 0&0&-1&0&0&0&0&0\\
 1&-1&0&0&0&0&0&0
 \end{pmatrix}.
\end{equation}
Direct evaluation gives $\det A=1$ and shows that $A$ maps the $19$
ray generators of $X_2$ bijectively onto those of the listed variety.
Since both fans are the fans over the faces of their smooth Fano
polytopes, this ray bijection identifies the fans as well.

There is also a finite verification of stability directly from the
database rays, independent of the path-family proof.  Its rational
base-polytope expression consists of $19$ bases with positive rational
coefficients summing to $1$.  The weighted sum of their incidence
vectors equals $(\delta_v)_v$, and their directed exchange graph is
strongly connected.  Lemma~\ref{lem:basecertificate} therefore proves stability.
These finite checks and the absence of any other stable record prove
uniqueness and maximality.
\end{proof}

\subsection{The varieties of Picard number \texorpdfstring{$n+2$}{n+2}}

For clarity, define the even-dimensional del Pezzo variety $V_{2m}$
as the smooth toric Fano variety whose fan is the fan over the faces
of the polytope
\begin{equation}\label{census:del-pezzo}
 \operatorname{conv}\{\pm e_1,\ldots,\pm e_{2m},
                  \pm(e_1+\cdots+e_{2m})\}\subset\ZZ^{2m}_\RR.
\end{equation}
This is the del Pezzo variety of Voskresenski\u\i{} and Klyachko
\cite{VoskresenskiiKlyachko1985}.  It has $4m+2$ rays and Picard number
$2m+2$; in particular $V_2=S_6$.  A toric fibre bundle with this fibre
means a Zariski locally trivial toric morphism with fibre $V_{2m}$,
equivalently a fan split by the base fan and the fan of $V_{2m}$ as in
Lemma~\ref{graph:split-collections}.

\begin{theorem}\label{census:boundary}
Suppose $4\le n\le7$, $X$ is a smooth toric Fano $n$-fold,
$\rho(X)=n+2$, and $T_X$ is stable.  Then either $n$ is even and
$X\cong V_n$, or $X$ is a toric fibre bundle with fibre $V_{2m}$,
$0<2m<n$, over a smooth projective toric variety $Y$ satisfying
$\rho(Y)=\dim Y$.  The numbers of isomorphism classes are as follows:
\[
 \begin{array}{c|rrrr}
 n&V_n&V_2\text{-bundles}&V_4\text{-bundles}&\text{total}\\ \hline
 4&1&2&0&3\\
 5&0&1&0&1\\
 6&1&7&2&10\\
 7&0&14&1&15.
 \end{array}
\]
\end{theorem}

\begin{proof}
For the $29$ stable varieties counted in
Theorem~\ref{census:lowdim}, we examine rational subspaces spanned by
rays in even dimension.  The fibre configurations are identified with
\eqref{census:del-pezzo} in integral coordinates.  In each of the $27$
proper-fibre cases, the non-fibre rays project bijectively onto the
rays of a smooth complete quotient fan.  Moreover, the maximal cones
of the original fan are exactly all sums of a fibre maximal cone and
a chosen lift of a quotient maximal cone.  Thus the full split-fan
condition holds, and the toric local triviality theorem applies.
The other two configurations are $V_4$ and $V_6$ themselves.

The base is projective: a base maximal cone lifts to a cone of the
original fan, and the orbit closure associated to any fixed fibre
maximal cone is a section isomorphic to the base.  Since $X$ is
projective, so is this closed subvariety.  Finally, counting fibre
and base rays in a split fan gives
\[
 \rho(X)=\rho(V_{2m})+\rho(Y),\qquad
 \dim X=2m+\dim Y.
\]
Substituting $\rho(X)=n+2$ and $\rho(V_{2m})=2m+2$ yields
$\rho(Y)=\dim Y$.
\end{proof}

The statement describes a finite part of the classification, not a
criterion asserting stability for every del Pezzo bundle of these
dimensions.  Nor does it imply that the base is a product of projective
lines.  Already in dimension six there are stable hexagon bundles over
other smooth toric bases of dimension and Picard number four.

\subsection{Reproducibility}

The ancillary package includes standalone rational-arithmetic programs
for the uniform path, end-attachment, branching and cycle
calculations. Their inputs are the rays, moment matrices, rank recurrences
and comparison bounds stated above. In particular, the end-attachment
and branching programs reproduce the finite minima, all-selected costs
and error bounds. The unequal-arm calculation evaluates all $500$
central choices for each of $7140$ mixed comparisons, including every
short-length case required by the contraction argument. The general-tree
theorem then uses an induction rather than an enumeration of tree shapes.
Its additional arithmetic consists of the limiting transition inequalities
and the explicit subdivision thresholds. From the ancillary
directory, the commands
\begin{quote}
\small\ttfamily\raggedright
python3 -S code/end\_attachment\_proof.py --uniform\\
python3 -S code/unequal\_spider\_proof.py --uniform\\
python3 -S code/two\_junction\_gluing.py\\
python3 -S code/subcubic\_thresholds.py\\
python3 -S code/research\_graph\_targets.py
\end{quote}
require only the Python standard library. The package retains the exact
rational outputs and $93$ positive basis decompositions for independently
checking selected finite examples. Running
\texttt{python3 -S code/verify\_certificate.py certificates/*.json}
checks all these decompositions. Numerical searches for decompositions
are not dependencies of their rational verification or of the uniform
stability proofs.

For the mixed-blowup construction, the tree criterion, the all-length
path theorem and the saturated-vertex comparison are proved by identities
and inequalities without a finite graph enumeration. The fourteen exact
fractions used in the star-transition proof are reproduced directly from
\eqref{branch:finite-sum} and its derivative by
\begin{quote}
\small\ttfamily\raggedright
python3 -S code/mixed\_star\_table.py --check-json\\
\hspace*{1em}data/mixed\_star\_table/signs.json
\end{quote}
The second line continues the same command. This calculation covers only
$m\le14$; the estimate \eqref{branch:star-tail} proves the assertion
for every $m\ge15$.

The finite rational inputs for the tree stability theorem and the local
degree-four and degree-five obstructions are reproduced by
\begin{quote}
\small\ttfamily\raggedright
python3 -S code/mixed\_tree\_arithmetic.py --check-json\\
\hspace*{1em}data/mixed\_tree\_arithmetic/signs.json
\end{quote}
This also checks the eighteen positive linear-functional values and six
adjacent pairs in Proposition~\ref{treemix:two-centres}. The conclusions
for arbitrary trees and connecting lengths use the monotonicity and
induction arguments in Section~\ref{sec:mixed-tree-threshold}.

The ancillary data distinguish input polytopes from varieties whose
stability has actually been decided.  The seven-dimensional records
cover all $1277$ varieties with $\rho\ge9$; the eight-dimensional records
cover precisely the $1582$ with $\rho\ge11$.  They retain identifiers,
ray coordinates, divisor degrees and the subspaces proving
non-stability.  The stable-boundary list is restricted to $\rho=n+2$
in dimensions four through seven.  In particular it does not include
$X_2$, whose Picard number is $n+3$.

For the stable eightfold, the rational expression in $19$ bases can be
checked using only the ray matrix, the normalized degrees and rational
arithmetic: verify each determinant is nonzero, each coefficient is
positive, their sum is $1$, the weighted incidence-vector identity
holds, and the directed exchange graph is strongly connected.  This check uses no fan
symmetry.  The identification matrix \eqref{census:identification}
gives a separate comparison with the graph construction.\footnote{\raggedright The
computations are implemented by \texttt{screen\_high\_rank.py},
\texttt{matroid\_certificate.py} and \texttt{verify\_certificate.py}.
The stable-eightfold expression is retained as
\texttt{P8D\_389596.json}.  Only completed calculations enter the
ancillary tables.}
